\section{Motivation and Bottleneck Analysis}
\label{sec:motivation}

Parallel decoding via diffusion is theoretically appealing, but fundamental bottlenecks prevent DLLMs from matching AR models in practice. We identify two categories of gaps---quality and systems---and analyze each below, summarizing the key evidence in Figure~\ref{fig:bottlenecks}.

\begin{figure*}[t]
\begin{center}
\begin{subfigure}[t]{0.32\textwidth}
\centering
\fbox{\rule[-.5cm]{0cm}{3.5cm} \rule[-.5cm]{0pt}{0cm}\hspace{3.8cm}}
\caption{\textbf{Infrastructure gap.} Per-component latency breakdown of AR (Qwen3-8B on SGLang) vs.\ SDAR vs.\ Strided DLLM (Ours) on H100. \yifan{@yuqing need to profile and generate this figure.}}
\label{fig:infra_bottleneck}
\end{subfigure}
\hfill
\begin{subfigure}[t]{0.32\textwidth}
\centering
\fbox{\rule[-.5cm]{0cm}{3.5cm} \rule[-.5cm]{0pt}{0cm}\hspace{3.8cm}}
\caption{\textbf{Compute economics.} Throughput vs.\ batch size for AR (Qwen3-8B on SGLang), SDAR, and Strided DLLM (Ours) on H100. AR scales; DLLM collapses at high concurrency; Ours matches AR scaling. \yifan{@yuqing need to benchmark and generate this figure.}}
\label{fig:compute_bottleneck}
\end{subfigure}
\hfill
\begin{subfigure}[t]{0.32\textwidth}
\centering
\fbox{\rule[-.5cm]{0cm}{3.5cm} \rule[-.5cm]{0pt}{0cm}\hspace{3.8cm}}
\caption{\textbf{Prefill--decode mismatch.} \textit{Placeholder:} Accuracy under matched vs.\ mismatched attention masks on representative benchmarks.}
\label{fig:prefill_bottleneck}
\end{subfigure}
\end{center}
\caption{\textbf{Bottleneck analysis.} Key gaps between DLLMs and AR models: (a)~AR serving infrastructure does not support diffusion-style denoising efficiently; (b)~DLLM parallel decoding consumes far more compute per token, collapsing throughput under concurrency; (c)~mismatched attention patterns between prefill and decoding degrade DLLM accuracy.}
\label{fig:bottlenecks}
\end{figure*}

\subsection{Quality Gap: Training Efficiency and Understanding--Generation Mismatch}
\label{sec:mot_quality}

\fan{This part is largely the background. the motivation should have our own experiments/figures.}

\paragraph{DLLM training is fundamentally less efficient.}
DLLMs consistently lag behind AR models in accuracy, and their training is fundamentally less efficient. Masked discrete diffusion models require up to $16\times$ more FLOPs to match the validation NLL of an AR equivalent~\citep{nie2024scalingmdlm}, widening to $64\times$ for continuous diffusion~\citep{gulrajani2023plaid}. As \citet{sdar2025} analyze, this stems from two sources: (1)~DLLMs optimize the evidence lower bound (ELBO), a loose upper bound of NLL that converges more slowly than cross-entropy~\citep{nie2024scalingmdlm, feng2025theoretical}; and (2)~masked training dilutes the supervision signal---with mask ratio $r$, only $r \cdot L$ of $L$ positions receive gradients per forward pass, while the remaining $(1{-}r) \cdot L$ positions consume the same compute but contribute nothing. At a typical $r{=}0.5$, each DLLM training step extracts only half the useful signal of an AR step. This asymmetry makes DLLM training less scalable: SDAR~\citep{sdar2025} requires 54B tokens of diffusion training yet still falls short of its base AR model. In contrast, our approach reaches \textbf{XX}\% of the base AR accuracy with only \textbf{10B} tokens (Table~\ref{tab:training_efficiency}), a $5.4\times$ reduction in token budget.
\yifan{Added our vs SDAR token budget comparison; need to fill in accuracy number.}
\fan{why this part belongs to quality gap? In fact, we can partition it into two parts. the high-level takeaway/background part (e.g., optimizing for ELBO) should be the leading part of this section, serving as the background while some details can be reallocated to quality and systems part.}

\paragraph{Attention mismatch between prefill and decoding degrades quality.}
A key advantage of AR models is that causal attention naturally unifies understanding and generation: the hidden state at position $i$ simultaneously encodes comprehension of all preceding context \emph{and} serves as the prediction signal for token $i{+}1$~\citep{wang2022zeroshot, fu2023decoderonly, behnamghader2024llm2vec}. DLLMs disrupt this unification by employing blockwise or bidirectional attention during generation, computing hidden states under a different pattern than the causal one used during prefill. To quantify this effect, we evaluate our Qwen3-8B-based DLLM under strict causal attention versus block-causal attention across representative benchmarks (Figure~\ref{fig:prefill_bottleneck}). Switching to block-causal decoding causes average accuracy to drop from \textbf{XX}\% to \textbf{XX}\%, a \textbf{XX}-point degradation. This is consistent with findings from \citet{nbdiff2025}, who report a similar drop from 48.6\% to 31.4\%. We provide further ablations in Section~\ref{sec:ablation}.
\yifan{Restructured: lead with argument, then our experiment numbers, cite NBDiff as support.}
\fan{change the writing so that 1) we introduce an argument (e.g., Block-causal causes accuracy drop); 2) introduce our experiment numbers; 3) (if needed) our results align with recent advances.}

\subsection{Systems Gap: Infrastructure Mismatch and Significant Compute Overhead}
\label{sec:mot_systems}

\paragraph{Existing infrastructure is built for AR, not diffusion.}
The modern LLM serving stack has been engineered around the sequential structure of AR decoding. Systems like vLLM~\citep{kwon2023vllm} and SGLang~\citep{zheng2024sglang} exploit the incremental nature of AR generation: each decode step appends exactly one token, enabling paged KV cache management, continuous batching, and fused attention kernels optimized for the append-one-token pattern. DLLMs break these assumptions---each denoising step processes a block of $N$ mask tokens simultaneously, requiring modifications to decoding kernels that permit only a single output token per step. The KV cache cannot be incrementally extended, and continuous batching becomes difficult because different requests may be at different denoising steps with different mask configurations. As a result, existing DLLM systems build custom inference pipelines from scratch, forgoing the extensive optimizations accumulated in AR serving stacks. To quantify this gap, we profile the per-component latency of AR (Qwen3-8B on SGLang), SDAR, and our Strided DLLM on the same hardware (Figure~\ref{fig:infra_bottleneck}). The breakdown reveals that SDAR spends \textbf{XX}\% of decode time on \textbf{XX}, compared to \textbf{XX}\% for AR. In contrast, our Strided DLLM achieves a latency profile of \textbf{XX}\%, closely matching AR---confirming that the infrastructure mismatch, not the model itself, is the dominant bottleneck, and that preserving causal structure eliminates it.
\yifan{@yuqing need to check the correctness and fill in numbers.}

\paragraph{Parallel decoding costs far more compute than it saves.}
To generate $N$ new tokens, an AR model performs $N$ forward passes, each processing a single new token position against the cached prefix, yielding a total cost proportional to $N$. A DLLM performs $T$ denoising steps over the same $N$ positions, plus an additional pass to commit the KV cache. The total cost is thus proportional to $(T{+}1) \times N$ token-positions---a $(T{+}1)\times$ computational overhead compared to AR. In latency-critical, single-request settings, this trade-off can be worthwhile~\citep{mercury2025, tidar2025}. However, in high-concurrency cloud environments where throughput under large batch sizes is the dominant objective, the extra compute makes DLLMs significantly less cost-effective than AR models. We measure throughput versus batch size for AR, SDAR, and our Strided DLLM (Figure~\ref{fig:compute_bottleneck}). AR throughput scales to \textbf{XX} tok/s at batch size 64, while SDAR collapses to \textbf{XX} tok/s. Our approach, by reusing the AR serving stack directly, achieves \textbf{XX} tok/s under the same conditions.
\yifan{@yuqing need to benchmark and fill in throughput numbers.}
